\subsection{Where this sits}\label{sec:related}

Conformal prediction has barely reached survey venues. \emph{Survey
Methodology} contains one conformal paper, \citet{wieczorek2023design}, which
gives design-based prediction for a unit drawn from the same finite population.
This journal contains two, both in one 2024 issue:
\citet{bersson2024optimal}, on small areas with area-specific frequentist
coverage, and \citet{michal2024model}, which is the closer antecedent and easy
to miss because its title does not name the method. That paper proposes a
scaled split-conformal procedure relaxing exchangeability for complex designs,
and documents that plain split conformal undercovers there. It differs from the
present setting in what carries the sampling error: its calibration units are
observations from a designed sample, whereas ours are \emph{estimated curves},
each with a design variance of its own, and it is that error in the calibration
objects rather than in the target that the correction studied here attempts to
remove.

The coverage-guarantee tradition this belongs to in this journal begins earlier,
with \citet{burris2020exact}, whose adaptive intervals attain exact area-level
coverage for every parameter value and are robust to misspecification of the
linking model. The competing line, second-order-correct empirical Bayes
intervals, culminating in \citet{yoshimori2014second}, is provably shorter than
direct intervals and is the benchmark any latent-target proposal must eventually
meet; we make no width claim against it, for the reason given in
Section~\ref{sec:discussion}.

On the estimand, complex-survey estimation of a distribution function is
well developed \citep{pasquazzi2016comparison,harms2006calibration}, but
simultaneous bands for one are not. The only \emph{Survey Methodology}
construction of simultaneous bands for a design-based estimated curve is
\citet{cardot2013confidence}, and its curve is a functional mean rather than a
distribution function. \citet{koffman2025function} is the closest existing
machinery to ours---a Rao--Wu--Yue--Beaumont bootstrap supporting joint bands
under a complex design---with a different estimand and no conformal component.

\paragraph{The resampling scheme, and a caution about it.} We use the
Rao--Wu--Yue rescaled bootstrap \citep{rao1992recent}, which
\citet{bessonneau2021with} treat for stratified multistage household surveys and
\citet{scarpa2026inference} report as outperforming alternatives in this journal
for such designs. One objection has to be met. \citet{saigo2007mean} recommends
an unrescaled, mean-adjusted scheme precisely \emph{because} no rescaling is
performed, on the grounds that rescaling sits awkwardly with estimating
distribution functions and quantiles. Our estimands are distribution functions,
so the objection applies. We answer it empirically rather than by assertion:
Section~\ref{sec:applications} reports every diagnostic under both schemes, and
the choice is consequential---the reliability gate opens in \VerMofmOpen{} of
\VerMofmN{} arms under the unrescaled scheme and \VerRwyOpen{} of \VerRwyN{}
under rescaling. We report the rescaled results as primary because the
unrescaled scheme understates between-unit variance in a stratified design, and
we report the difference rather than absorbing it.
